\scanpage{51}
\mspage{5}
en notant que pour deux droites vectorielles en dualité $F,F'$, on a $F^*\simeq F'^*$ can., --- mais cet isom. (20) \inserted{entre $(P\Sigma)_x$ et $(F\mathbin\wedge\omega)^*$} \deleted{\unsure{doit}} \uline{renverse} l'\uline{action} de $\mathbb C^*$.

On trouve ainsi des iso
\[
(21)\quad\left\{\begin{aligned}P\Sigma&\simeq(V(\omega))^*\\P\Sigma&\simeq V^*\end{aligned}\right.
\]
\diagram{\begin{tikzcd}[ampersand replacement=\&, column sep=small]
\deleted{\unsure{F}}\otimes\deleted{\unsure{F'}}\arrow[r] \&\deleted{\unsure{D}}\arrow[r]\&\deleted{\unsure{\mathbb C}}\\
\deleted{\unsure{F/e}}\&\&
\end{tikzcd}}
\margin{qui transforme le couple $(x,e\in(P\Sigma)_x\simeq\check F(\omega)\simeq F^\perp)$ \unsure{en} $F(\omega)^*\subset V(\omega)^*$ ; soit en $\check e$, considéré comme élém\unsure{$^t$} de $(F^\perp)^*\subset V^*$ soit en $e$, considéré comme élém\unsure{$^t$} de $\mathbb C^*$. Ces iso (21) \unsure{sont compatibles} à l'action de $D$ (\unsure{le groupe de structure}).}

Mais on a un isom canonique (provenant du produit extérieur
\[
V\simeq\check V(\omega)\qquad\deleted{et}
\]
et un isom canonique (provenant de la forme quadratique
\[
V\longrightarrow\check V\qquad\text{donc}\qquad\check V(\omega)\simeq V(\omega)
\]
d'où un isom canonique
\[
(21)\qquad V\xrightarrow{\sim}V(\omega)
\]
qui explicite le passage de l'un à l'autre isom. \deleted{\unsure{commutant d'ailleurs aux actions de $\mathbb C^*$}} (21). D'ailleurs, (21) induit \deleted{$D$, et par \unsure{transposant} l'identité} un isom
\diagram{\begin{tikzcd}[ampersand replacement=\&]
\deleted{\Sigma\simeq P(V)}\arrow[r,"\deleted{\mathrm{id}}"]\&\deleted{\Sigma\simeq P(V(\omega))}
\end{tikzcd}}
\diagram{\begin{tikzcd}[ampersand replacement=\&]
\text{(22)}\quad\Sigma=P(V)\arrow[r,"\sigma"]\&P(V(\omega))\quad(\deleted{\unsure{\Sigma}}\simeq\Sigma=P(V))
\end{tikzcd}}
(\deleted{qui est} la \deleted{\unsure{rotat\textsuperscript{n}}} symétrie p.r. à l'axe polaire, \inserted{qui du point de vue linéaire s'exprime} \deleted{i.e. l'\unsure{échange} de} comme le passage de $F$ à $F^\perp$). \deleted{et}

\deleted{par cet iso., \unsure{\illegible{1}} à $F$ correspond : $F'(\omega)(\simeq F^\perp)$ dans $V(\omega)$ \unsure{avec} \inserted{(où $F'\subset V$, en fait $F'=F^\perp$)} \unsure{on a} $F\longmapsto F'(\omega)$ i.e. $F(\omega)\simeq F'$.}
\mspage{50}
